% ===========================================================================
%  conclusion.tex  --  Заключение
% ===========================================================================

\newpage
\section*{Заключение}
\addcontentsline{toc}{section}{Заключение}

Настоящая работа была посвящена исследованию условий, при которых
обучаемая декомпозиция временных рядов на~компоненты и~нейросетевые
механизмы учета долгосрочных зависимостей дают практическое
преимущество по~сравнению с~фиксированными статистическими подходами,
а~также ситуаций, в~которых более простые методы оказываются
предпочтительнее. Для ответа на~этот вопрос был разработан
масштабный экспериментальный контур, включающий 11~моделей
(4~классические, 7~нейросетевых) и~5829~временных рядов
из~наборов M3 Competition (2829~рядов: 645~yearly, 756~quarterly,
1428~monthly) и~M4 Competition (3000~рядов: 1000~quarterly,
1000~monthly, 1000~daily). Эксперимент охватывает прямое
прогнозирование, итеративный режим и~предобучение на~корпусе рядов.

Основные результаты работы можно сформулировать следующим образом.

\textbf{1. Классические модели лидируют в~per-series режиме,
но нейросетевые конкурентоспособны на~длинных рядах.}
На~полном M3 Competition (2829~рядов) Seasonal Naive (15.74\,\%),
ETS (16.04\,\%) и~Auto-ARIMA (16.14\,\%) лидируют по~среднему
взвешенному sMAPE. TimesNet (17.65\,\%) --- лучшая нейросетевая
модель, отставая от~Seasonal Naive менее чем на~2~п.п. На~M4
Competition (1000~рядов на~категорию) картина аналогична:
Auto-ARIMA (12.51\,\%) и~ETS (12.61\,\%) лидируют, TimesNet
(14.29\,\%) ближе всего к~классическим. Однако на~ежедневных
рядах M4 (длина до~9919 наблюдений) TimesNet (3.62\,\%) уступает
Auto-ARIMA (3.22\,\%) лишь на~0.4~п.п., а~PatchTST (4.44\,\%),
N-BEATS (4.46\,\%) и~Helformer (4.49\,\%) демонстрируют
сопоставимые результаты.

\textbf{2. Длина ряда является критическим фактором для
нейросетевых моделей.} Экспериментально установлена граница порядка
50--100 наблюдений, ниже которой нейросетевые модели систематически
уступают статистическим, а~выше --- становятся конкурентоспособными.
На~месячных рядах M3 ($\sim$126 наблюдений) Helformer (17.98\,\%)
и~N-BEATS (18.26\,\%) сопоставимы с~ETS (17.42\,\%). На~ежедневных
рядах M4 (до~9919 наблюдений) все нейросетевые модели
конкурентоспособны с~классическими.

\textbf{3. Предобучение кардинально повышает качество нейросетевых
моделей.} Масштабный эксперимент с~предобучением на~шести частотных
категориях (M3 yearly/quarterly/monthly, M4 quarterly/monthly/daily)
и~семи моделях выявил дифференцированный эффект переноса знаний.
DLinear показывает наибольший средний выигрыш (+31.5\,\%),
Autoformer~--- +25.5\,\%, Helformer~--- +18.2\,\%.
На~M4~daily PatchTST с~предобучением (2.36\,\%) превосходит
Auto-ARIMA (3.22\,\%) и~ETS (3.46\,\%), демонстрируя, что при
наличии корпуса однородных рядов нейросетевые модели способны
превзойти классические. TimesNet, напротив, деградирует при
предобучении на~коротких рядах ($-$52.6\,\% на~M3~yearly),
но улучшается на~длинных (M4~daily: +13.5\,\%), что свидетельствует
о~пороговой зависимости эффекта предобучения от~длины ряда.

\textbf{4. Обучаемая декомпозиция эффективна при достаточном
объёме данных.} Модели с~обучаемой декомпозицией (TimesNet, N-BEATS)
демонстрируют переход от~неэффективной к~конкурентоспособной
при увеличении длины ряда. TimesNet с~двумерной декомпозицией
(средний ранг~4.17) входит в~четвёрку лидеров на~пяти из~шести
категорий, уступая лишь классическим моделям. Трансформерные
модели с~обучаемой декомпозицией серийного типа (Autoformer,
FEDformer) оказались наименее эффективными в~per-series режиме
(средний ранг 9.33 и~10.33), однако при предобучении значительно
улучшаются.

\textbf{5. Итеративное прогнозирование выявляет хрупкость обучаемой
декомпозиции.} При переходе от~прямого к~итеративному
прогнозированию модели с~обучаемой декомпозицией демонстрируют
наибольшую деградацию: FEDformer $+5.01$~п.п., TimesNet
$+5.16$~п.п. PatchTST ($+2.18$~п.п.), работающий с~патчами
без явной декомпозиции, оказался наиболее устойчивым трансформером.
Классические модели (ETS, Auto-ARIMA) практически нечувствительны
к~режиму прогнозирования.

На~основании полученных результатов можно сформулировать следующие
\textbf{практические рекомендации}:

\begin{itemize}[nosep]
\item для одномерных рядов с~выраженной сезонностью при малом объёме
  данных ($<$100 наблюдений) ETS и~Auto-ARIMA являются лучшим
  выбором;
\item при длине ряда свыше 100 наблюдений целесообразно рассмотреть
  TimesNet или N-BEATS как альтернативу классическим моделям,
  особенно при наличии нелинейных паттернов;
\item при наличии корпуса однородных рядов предобучение нейросетевых
  моделей (особенно DLinear, Autoformer, Helformer) даёт
  значительное улучшение (+18--32\,\%), что делает их
  конкурентоспособными и~превосходящими классические подходы;
\item для задач, требующих итеративного прогнозирования, следует
  отдавать предпочтение моделям без сложной обучаемой декомпозиции
  (PatchTST, DLinear) или классическим подходам;
\item тип декомпозиции должен соответствовать режиму обучения:
  фиксированная декомпозиция (ETS, DLinear) наиболее устойчива,
  обучаемая (TimesNet, Autoformer) --- наиболее адаптивна при
  достаточном объёме данных.
\end{itemize}

Таким образом, проведённое исследование показывает, что ответ
на~вопрос о~превосходстве обучаемой декомпозиции не~является
бинарным: эффективность нейросетевых механизмов критически зависит
от~длины ряда, режима обучения и~доступного объёма данных.
Предобучение на~корпусе рядов позволяет эффективно преодолеть
порог минимальной длины и~вывести нейросетевые модели на~уровень,
превосходящий классические подходы.
